\begin{figure}[htbp]
\centering
\begin{tikzpicture}
\begin{loglogaxis}[
    labbook standard,
    xlabel={Klein Perturbation Parameter $\epsilon$},
    ylabel={Boundary Energy Splitting $|\delta E_{\mathrm{site}}(\boldsymbol{\epsilon})|$},
    legend pos=north west,
    xmin=1e-3, xmax=0.06,
    ymin=1e-9, ymax=1e-4
]
    \addplot[labbookaccent, thick, mark=*, mark size=2.2pt]
        table[x index=0, y index=4] {../../data/processed/benchmarks/klein_duality_splitting.dat};
    \addlegendentry{Numerical Finite-Cell Splitting}

    % Cubic reference line
    \addplot[labbookcrimson, dashed, thick, domain=1.5e-3:0.05] {0.38 * x^3};
    \addlegendentry{Odd-Parity Cubic $\mathcal{O}(\epsilon^3)$ Law}
\end{loglogaxis}
\end{tikzpicture}
\caption{Four-sublattice Klein duality finite-cell boundary splitting. Halving the perturbation step size decreases the energy difference by a factor of 8 ($2^3$), confirming the proven odd-parity cubic scaling law $\delta E = c_3 \epsilon^3 + \mathcal{O}(\epsilon^5)$.}
\label{fig:m1_4_5_klein_splitting}
\end{figure}
